\documentclass[11pt,letterpaper]{article}
\usepackage[margin=0.8in]{geometry}
\usepackage[T1]{fontenc}
\usepackage{lmodern,microtype,amsmath,amssymb,booktabs,tabularx,graphicx,caption}
\usepackage[hidelinks]{hyperref}
\setlength{\parindent}{0pt}\setlength{\parskip}{5pt}
\captionsetup{font=small,labelfont=bf,skip=4pt}
\setlength{\textfloatsep}{8pt}\setlength{\floatsep}{7pt}
\setlength{\emergencystretch}{2em}
\graphicspath{{figures/}}
\newcommand{\CartBasicMean}{500.000}
\newcommand{\CartBasicSD}{0.000}
\newcommand{\CartBasicClipped}{0}
\newcommand{\PongBasicMean}{-18.610}
\newcommand{\PongBasicSD}{1.681}
\newcommand{\PongBasicClipped}{1542}
\newcommand{\CartBaselineMean}{260.308}
\newcommand{\CartBaselineSD}{57.013}
\newcommand{\CartBaselineClipped}{0}
\newcommand{\PongBaselineMean}{-18.446}
\newcommand{\PongBaselineSD}{1.851}
\newcommand{\PongBaselineClipped}{466}
\newcommand{\CartVarianceReduction}{84.22}
\newcommand{\PongInitialVarianceReduction}{94.38}
\newcommand{\PongReferenceDiscounted}{-1.309180}
\newcommand{\CartSeedRows}{%
No baseline & 1668 & 500.000 & 0.000 & 1.000 \\
No baseline & 1669 & 271.824 & 28.856 & 0.000 \\
No baseline & 1670 & 483.886 & 65.298 & 0.928 \\
Learned baseline & 1668 & 260.308 & 57.013 & 0.008 \\
Learned baseline & 1669 & 360.516 & 79.141 & 0.100 \\
Learned baseline & 1670 & 349.940 & 76.321 & 0.092 \\%
}
\newcommand{\CartTraceRows}{%
0 & 0.01276 & -0.01188 & -0.04659 & -0.03077 & 0.1019 & 0 \\
1 & 0.01252 & -0.20630 & -0.04721 & 0.24686 & 0.8335 & 1 \\
2 & 0.00840 & -0.01054 & -0.04227 & -0.06033 & 0.0700 & 0 \\
3 & 0.00818 & -0.20503 & -0.04348 & 0.21872 & 0.7711 & 0 \\
4 & 0.00408 & -0.39951 & -0.03910 & 0.49738 & 0.9930 & 1 \\
5 & -0.00391 & -0.20385 & -0.02915 & 0.19263 & 0.7348 & 1 \\
6 & -0.00798 & -0.00833 & -0.02530 & -0.10911 & 0.0428 & 0 \\
7 & -0.00815 & -0.20308 & -0.02748 & 0.17549 & 0.6828 & 1 \\
8 & -0.01221 & -0.00757 & -0.02397 & -0.12574 & 0.0337 & 0 \\
9 & -0.01236 & -0.20234 & -0.02649 & 0.15929 & 0.6266 & 0 \\%
}
\newcommand{\PongVarianceRows}{%
EMA-baseline actor & 4096 & 381.66 & 333.33 & 0.873 \\
EMA-baseline actor & 8192 & 160.81 & 163.52 & 1.017 \\
EMA-baseline actor & 12288 & 217.75 & 201.19 & 0.924 \\
EMA-baseline actor & 16384 & 118.86 & 120.18 & 1.011 \\
Basic actor & 4096 & 437.78 & 53.44 & 0.122 \\
Basic actor & 8192 & 434.45 & 291.59 & 0.671 \\
Basic actor & 12288 & 193.44 & 199.18 & 1.030 \\
Basic actor & 16384 & 200.39 & 224.50 & 1.120 \\%
}
\newcommand{\PongTraceRows}{%
0 & 0.2772 & 2 & 0 & 0.59704 & 1 \\
1 & 0.0007 & 3 & 0 & 0.60307 & 0.99 \\
2 & 0.0064 & 3 & 0 & 0.60917 & 0.98 \\
3 & 0.0003 & 3 & 0 & 0.61532 & 0.97 \\
4 & 0.9862 & 2 & 0 & 0.62154 & 0.961 \\
5 & 0.0000 & 3 & 0 & 0.62781 & 0.951 \\
6 & 0.1081 & 3 & 0 & 0.63415 & 0.941 \\
7 & 1.0000 & 2 & 0 & 0.64056 & 0.932 \\
8 & 0.6196 & 2 & 0 & 0.64703 & 0.923 \\
9 & 0.0000 & 3 & 0 & 0.65357 & 0.914 \\
93 & 0.8358 & 2 & 1 & 1.52031 & 0.393 \\
1797 & 0.0007 & 3 & -1 & -1.00000 & 1.43e-08 \\%
}
\newcommand{\CartRateRows}{%
0.01 & 418.570 & 323.588 \\
0.005 & 163.178 & 357.661 \\
0.002 & 247.102 & 264.481 \\%
}
\newcommand{\CartControlRows}{%
Always left & 9.34 & 9.50 & 7.75 \\
Always right & 9.38 & 9.50 & 7.78 \\
Basic, greedy & 469.77 & 488.44 & 8.70 \\
Baseline, greedy & 283.61 & 284.77 & 13.33 \\%
}
\newcommand{\GithubUsername}{FailedMath3U}
\newcommand{\RepositoryName}{MLDS}
\newcommand{\RepositoryLink}{\url{https://github.com/FailedMath3U/MLDS}}
\newcommand{\SubmissionNotice}{}

\newcommand{\resultplots}[6]{%
\begin{figure}[!ht]\centering
\includegraphics[width=.85\linewidth,height=.27\textheight,keepaspectratio]{#1}
\caption{#2: every training episode reward and trailing 100-episode mean. #3}
\end{figure}
\begin{figure}[!ht]\centering
\includegraphics[width=.85\linewidth,height=.25\textheight,keepaspectratio]{#4}
\caption{#2: 500 final-policy evaluation rewards. Mean #5; population SD #6.}
\end{figure}}
\begin{document}
\begin{center}{\Large\bfseries Assignment 1: Policy Gradient}\\[3pt]
CartPole-v1 and Pong-v5, Parts 1 and 2\end{center}
\textbf{GitHub username:} \GithubUsername\quad\textbf{Repository:} \RepositoryName\\
\RepositoryLink

\section{Algorithm and experimental protocol}
The PyTorch implementation uses REINFORCE: collect complete episodes, compute discounted reward-to-go, then update the policy once per batch. The policy $\pi_\theta(a_t\mid o_t)$ gives action probabilities from network input $o_t$; $\theta$ denotes all actor weights and biases. At decision $t$, the sampled action is $a_t$ and the subsequent reward is $r_t$. For episode length $T$ and discount $\gamma$, optimize
\[
J(\theta)=\mathbb E_{\pi_\theta}\!\left[\sum_{t=0}^{T-1}\gamma^t r_t\right],
\qquad G_t=\sum_{k=t}^{T-1}\gamma^{k-t}r_k.
\]
$G_t$ is the discounted reward remaining from decision $t$, calculated backward with zero continuation at the episode end. For $B$ episodes, Adam minimizes
\[
\mathcal L=-\frac{c}{B}\sum_{i=1}^{B}\sum_{t=0}^{T_i-1}
\gamma^t\log\pi_\theta(a_t^{(i)}\mid o_t^{(i)})(G_t^{(i)}-b_t^{(i)}).
\]
Here $i$ indexes episodes, $b_t$ is the baseline (zero in Part 1), and $c$ is a fixed positive scale. The code sums decision contributions and divides by the number of episodes, not their total length. Returns and baseline predictions are held fixed during actor differentiation. The two output logits define a categorical action distribution. There is no return normalization or centering. Pong accumulates gradients in chunks but takes only one optimizer step per complete batch.

\begin{center}\small
\begin{tabular}{lrr}\toprule
Actor setting & CartPole & Pong\\\midrule
Network & $4\to64$ tanh $\to2$ logits & $6400\to200$ ReLU $\to2$ logits\\
Discount $\gamma$; loss scale $c$ & 0.95; 0.05 & 0.99; 1\\
Adam learning rate & 0.01 & 0.001\\
Episodes per update; total per seed & 32; 8,000 & 8; 20,000\\
Training seeds & 1668, 1669, 1670 & 1668\\
Gradient-norm cap & 10 & 5\\
Hardware & CPU & One GPU\\\bottomrule
\end{tabular}\end{center}

\textbf{Evaluation protocol.} Each required histogram uses 500 fresh episodes from the final fixed-budget policy, with no learning. Actions are sampled, not chosen greedily. Test reset seeds are separate from training and shared across methods. CartPole seed 1668 was fixed before the full runs for the main plots; all three seeds are reported. Every figure uses undiscounted episode reward $\sum_t r_t$; SD is the population standard deviation across its 500 evaluations.

\textbf{Alternative Pong learning rate.} On identical real-game data, one update at 0.001 produced mean action-distribution KL divergence $1.996\times10^{-6}$, versus $3.423\times10^{-4}$ at 0.01. This motivated a smaller step; it does not prove greater long-run stability. Clipping affected 61.68\% of basic Pong updates and 18.64\% with a baseline. No recorded CartPole update reached its cap.

\clearpage
\section{Part 1(a): CartPole-v1 without a baseline}
\textbf{Model.} The four inputs are cart position, cart velocity, pole angle and pole angular velocity. Actions 0 and 1 apply left and right force. The default reward is $+1$ per transition, including failure. Physical failure occurs beyond cart position $\pm2.4$ or pole angle $\pm12^\circ$; a time limit stops episodes at 500 steps. Each reset independently samples the four components within $(-0.05,0.05)$ using recorded episode seeds. We optimize this capped episode, so returns stop at either end flag without bootstrapping. The critic in Part 2 additionally observes time; the actor does not.

\textbf{Result.} The primary final policy scored \textbf{\CartBasicMean{} $\pm$ \CartBasicSD}. All 500 evaluations reached the limit without physical failure. An untrained network scored $20.926\pm11.467$ and a uniform random policy $22.782\pm12.058$ over 500 episodes each. Learning was substantial but not stable throughout training; Section~\ref{sec:tuning} tests the large dips directly.
\resultplots{cartpole_no_baseline_training.pdf}{CartPole, Part 1}{Seed 1668; 8,000 complete episodes.}{cartpole_no_baseline_hist.pdf}{\CartBasicMean}{\CartBasicSD}

\clearpage
\section{Part 1(b): Pong-v5 without a baseline}
\textbf{Model.} We use \texttt{ALE/Pong-v5}, the namespaced Pong-v5 environment, with only actions 2/3 (ALE labels RIGHT/LEFT; vertical paddle movement). Rewards are $+1$ when the agent scores, $-1$ when the opponent scores, and zero otherwise. Episodes are complete games, not points. The assignment's crop, factor-two downsampling and background removal produce binary $80\times80$ frames. The network receives the flattened signed difference of consecutive frames; the first input after reset is zero.

This is a reactive policy in a partially observable process: $o_t$ supplies motion cues but is not the full emulator state $s_t$. Unchanged image regions cancel, and the network has no recurrent memory. Transition probabilities are supplied implicitly by the emulator, not learned. Defaults retain four emulator frames per decision, sticky-action probability 0.25 and a 108,000-frame cap.

\textbf{Result.} The final policy scored \textbf{\PongBasicMean{} $\pm$ \PongBasicSD}, versus an untrained reference of $-20.250\pm0.935$ over 32 different games. All 500 test games were losses (range $-21$ to $-10$), with no truncations. \textbf{The policy did not learn to win.} The reference comparison is descriptive, not a paired test.
\resultplots{pong_no_baseline_training.pdf}{Pong, Part 1}{Seed 1668; 20,000 full games.}{pong_no_baseline_hist.pdf}{\PongBasicMean}{\PongBasicSD}

\clearpage
\section{Part 2(a): CartPole with a learned baseline}
\textbf{Choice and implementation.} A separate network with inputs $(o_t,t/500)$, 64 tanh units and one output predicts $G_t$. Time matters because fewer rewards remain near the episode limit. Before each actor update, its predictions are computed without gradients. Afterward, five Adam steps at rate 0.01 fit the observed Monte Carlo returns by mean squared error. There is no temporal-difference bootstrap. Actor settings, initial actor weights and episode budgets match Part 1.

Subtracting an action-independent baseline preserves the expected \emph{unclipped} policy gradient:
\[
\mathbb E_{a\sim\pi_\theta}[b\nabla_\theta\log\pi_\theta(a\mid o)]
=b\nabla_\theta\sum_a\pi_\theta(a\mid o)=0.
\]
Here $\nabla_\theta$ is the vector of derivatives with respect to actor parameters. This identity does not guarantee lower variance for an inaccurate baseline.

\textbf{Result.} The final model scored \textbf{\CartBaselineMean{} $\pm$ \CartBaselineSD}: four episodes reached 500, and 496 ended in physical failure. It improved over random control but underperformed the primary basic policy. The variance experiment in Section~\ref{sec:cartvalidation} distinguishes a baseline's local variance benefit from successful training.
\resultplots{cartpole_baseline_training.pdf}{CartPole, Part 2}{Seed 1668; 8,000 complete episodes.}{cartpole_baseline_hist.pdf}{\CartBaselineMean}{\CartBaselineSD}

\clearpage
\section{Part 2(b): Pong with a constant baseline}
\textbf{Choice and implementation.} The baseline is a scalar exponential moving average (EMA), held constant within each batch. Starting at $b_0=0$, update $k$ uses only the previous value; afterward,
\[
\bar G_k=\frac{\sum_{i,t}\gamma^tG_t^{(i)}}{\sum_{i,t}\gamma^t},
\qquad b_{k+1}=0.95b_k+0.05\bar G_k.
\]
$k$ counts actor updates. Thus the current sampled actions do not determine the baseline used for their own gradient. Apart from this subtraction, the actor architecture, initial weights, environment, budgets and action restrictions match Part 1.

\textbf{Result.} The final model scored \textbf{\PongBaselineMean{} $\pm$ \PongBaselineSD}, 0.164 above the basic method. All 500 games were losses (range $-21$ to $-5$), with no truncations. One training seed does not establish a general benefit. Later paired probes found no consistent variance reduction from the actual EMA (Section~\ref{sec:pongvalidation}). Frequent gradient clipping also means actual updates differ from the unclipped estimator used in the baseline identity.
\resultplots{pong_baseline_training.pdf}{Pong, Part 2}{Seed 1668; 20,000 full games.}{pong_baseline_hist.pdf}{\PongBaselineMean}{\PongBaselineSD}

\clearpage
\section{CartPole: edge cases and validation}\label{sec:cartvalidation}
\textbf{Boundary crossings and unavoidable failure.} Supplementary checks tested both actions at both signed cart/pole limits, and just inside each limit with outward velocity: all 16 predictions matched. Equality at the threshold is allowed; crossing it terminates. For example, $x=2.39$ and $\dot x=1$ give $x'=2.41$ under either action because the default Euler step uses $x'=x+0.02\dot x$. Force changes velocity too late to avoid that transition's failure. The terminal reward remained $+1$.

\textbf{Failure and time limit together.} Setting a one-step limit and forcing a boundary crossing returned both \texttt{terminated=True} and \texttt{truncated=True}, with one reward of 1. A score of 500 alone therefore cannot identify failure-free survival; both flags are retained. The separate original time-limit check verified truncation without failure.

\textbf{Different episode lengths.} The actual collector was exercised with controlled episodes of lengths 1, 1 and 3. It recorded no post-end transitions; the one-step return was exactly 1. The combined loss gradient agreed with the mean of separately computed episode gradients to a maximum component error of $1.863\times10^{-9}$. This checks trajectory separation and averaging by episode rather than introducing length-dependent normalization.

\textbf{Negative controls and recovery limits.} Frozen primary actors and fixed-action controls were evaluated on 64 matched starts per condition, without retraining. Trained actors used \emph{greedy} actions here, separately from the required stochastic evaluations:
\begin{center}\small
\begin{tabular}{lrrr}\toprule
Policy & Default resets & Tilted, moving outward & Near track edge\\\midrule
\CartControlRows\bottomrule
\end{tabular}\end{center}
Entries are mean episode rewards. Tilted starts set pole angle to $\pm0.15$ radians and angular velocity to $\pm0.5$; edge starts set cart position to $\pm2.3$ and velocity to $\pm0.5$. Signs alternate; other components retain seeded reset values. Fixed-direction controllers fail quickly, while trained policies balance from default and tilted starts. Both struggle near the track edge. These exploratory stress tests expose limits of standard-reset results; they do not establish that every edge start is recoverable. The basic actor's greedy default score also shows why evaluation action selection must be specified.

\textbf{Baseline variance.} At a frozen initial actor, a separate critic fitted on 512 trajectories was measured on 512 different, paired trajectories. Summed sample variances of gradient components fell from 5.338262 to 0.842421 (84.22\%). The mean-gradient difference was 0.0474 versus paired standard-error magnitude 0.1176. This supports variance reduction at that actor, but the critic received 150 fitting steps; it does not establish the same benefit for five-step online fitting throughout training.

\textbf{Decision inspection.} The saved ten-step trace contains all four observations, probabilities, sampled actions and rewards. At step 3, right had probability 0.7711 but left was sampled; the pole subsequently remained within bounds. This is compatible with stochastic action selection, not evidence of a reversed action mapping. A separate fixed-seed, ten-action replay reproduced identical state sequences.

\clearpage
\section{Pong and shared gradient validation}\label{sec:pongvalidation}
\textbf{Scoring is not an episode end.} For rewards $(0,1,0,-1)$ at $\gamma=0.99$, independently calculated returns are
\[
(G_0,G_1,G_2,G_3)=(0.019701,0.0199,-0.99,-1).
\]
The implementation matched within $10^{-6}$. In particular, $G_1=1-0.99^2$ includes the later loss; incorrectly restarting returns at the positive point would give 1. Complete-game collection was checked separately. Both end flags stop collection, and frame history is reset between games.

\textbf{Serving under restricted actions.} Eight complete preflight games using only actions 2/3 each produced 21--22 scoring events and reached natural termination. No startup FIRE/NOOP wrapper was added. This verifies progress under the assigned restriction, not that FIRE has no effect in unrestricted play.

\textbf{Sticky actions.} The emulator can repeat its previously executed action with probability 0.25 per emulator frame. Requested and executed actions therefore need not coincide throughout a four-frame decision. We retain this stochasticity and test seeded replay, rather than assuming deterministic action execution: identical reset seeds and requested-action sequences reproduced frame/reward hashes. This establishes reproducibility, not robustness across all seeds.

\textbf{Visual input.} Known-pixel images verified the crop, factor-two sampling, background removal and foreground preservation. Subtraction uses signed values, so disappearing pixels become $-1$ rather than unsigned overflow. Real differences were in $\{-1,0,1\}$; reset input was zero. The saved game trace includes the first ten decisions, a point and the ending, with corresponding image differences. Its positive point at decision 93 and final loss at 1,797 illustrate how rewards can be far apart within one episode.

\textbf{Does the Pong baseline reduce variance?} An independently fitted candidate constant reduced summed gradient variance from 369.320 to 20.763 on 16 held-out games (94.38\%). It was fitted on eight other games and was \emph{not} the initial training EMA, which starts at zero. Later probes used the actual frozen EMA and 16 paired games at the same actor. Variance ratios (with/without baseline) at 4,096, 8,192, 12,288 and 16,384 games were 0.873, 1.017, 0.924 and 1.011. Reduction was not sustained. Small probe samples and the absence of ratio confidence intervals limit precision; all probes measure unclipped gradients.

\textbf{Analytic gradient controls for both parts.} In a two-step task with Bernoulli action probability $p=\operatorname{sigmoid}(\theta)$ and reward equal to the action, the exact objective is $(1+\gamma)p$. Enumeration over all four action sequences reproduced its derivative $(1+\gamma)p(1-p)$, with and without a constant baseline: error $8.52\times10^{-9}$ for CartPole's discount and below $10^{-7}$ for Pong's. Setting discount to zero recovered immediate rewards. Zero Pong returns produced zero basic-policy gradient. A one-step actor check with target 1 changed the selected action's probability from 0.5 to \PositiveAdvantageProbability{} with baseline 0, and to \NegativeAdvantageProbability{} with baseline 2. This verifies both signs of the advantage, rather than assuming positive game rewards require every action probability to increase.

These controls test the estimator and task handling. They do not establish that the Pong representation, optimizer settings or training budget were adequate to learn winning play.

\clearpage
\section{Tuning and interpretation}\label{sec:tuning}
\textbf{All CartPole seeds.} Each row is a final policy evaluated on 500 stochastic episodes at the original actor rate 0.01:
\begin{center}\small
\begin{tabular}{lrrrr}\toprule
Method & Seed & Mean & Population SD & Fraction at 500\\\midrule
\CartSeedRows\bottomrule
\end{tabular}\end{center}
The baseline improves seed 1669, but not the other two. The following additional runs changed only actor rate, keeping seeds, 8,000-episode budgets and 500 evaluations per model; critic rate remained 0.01:
\begin{center}\small
\begin{tabular}{rrr}\toprule
Actor rate & Basic: mean across three seeds & Baseline: mean across three seeds\\\midrule
\CartRateRows\bottomrule
\end{tabular}\end{center}
Smaller rates did not consistently improve final reward. These exploratory comparisons do not identify an optimal rate or establish that a baseline generally helps.

\textbf{Are the training dips just episode noise?} Both seed-1668 runs were replayed exactly: every original training reward and final parameter hash matched. Before/during checkpoints selected from those curves were then frozen and evaluated on 500 fresh paired reset/action seeds:
\begin{center}\small
\begin{tabular}{lrrr}\toprule
Method & Before mean & During mean & Change [paired 95\% interval]\\\midrule
Basic (5,120 $\to$ 6,400) & 453.108 & 83.114 & $-369.994\;[-377.430,-362.038]$\\
Baseline (3,072 $\to$ 6,144) & 440.466 & 215.686 & $-224.780\;[-231.886,-217.436]$\\\bottomrule
\end{tabular}\end{center}
Intervals use 5,000 paired bootstrap replicates. A further 2,000 episodes with each earlier actor frozen gave 100-episode mean ranges of 430.11--472.29 (basic) and 413.00--453.82 (baseline); continued training reached 74.64 and 211.37. Thus the selected dips reflect policy changes beyond measured fixed-policy episode noise. Noisy updates may still cause those changes. Selection was exploratory; intervals condition on these actors and do not measure uncertainty across training seeds. Rolling windows overlap, so their ranges are descriptive.

\textbf{Objective versus displayed score.} For CartPole, discounted return is $(1-0.95^T)/0.05$, close to 20 long before 500 steps. The primary baseline's mean discounted return was 19.999743 despite its lower survival time. In Pong, mean discounted return improved from $-1.309180$ for the untrained reference to 0.313913 (basic) and 0.293346 (baseline), while all final games were losses. A reward 1,000 decisions away has weight $0.99^{1000}\approx4.32\times10^{-5}$. This suggests a credit-assignment difficulty; it does not identify the cause of failure. Partial observation, optimization and the single training seed remain unresolved limitations.

\clearpage
\section{Execution and evidence}
Paths are relative to the repository's \texttt{assignment1/}. Retrain in a fresh copy to preserve saved runs. Versions: Python 3.12, PyTorch 2.14.0, NumPy 2.5.3, Gymnasium 1.2.2; Pong also uses ALE 0.11.2.

\textbf{CartPole:}
\begin{verbatim}
cd mlds_assignment1_cartpole
python -m pip install -r requirements.txt
python cartpole.py suite --episodes 8000 --lr 0.01
python validate.py
python validate_task_edges.py
python instability_experiment.py
python plot_instability_experiment.py
\end{verbatim}
The suite trains/evaluates six actors. \texttt{validate.py} writes analytic-gradient, return and baseline-variance evidence to \texttt{validation/}. \texttt{validate\_task\_edges.py} writes boundary, collector, advantage-sign and recovery evidence to \texttt{task\_validation/}. The instability scripts reproduce replay, frozen-policy comparisons and bootstrap results. Each output folder contains numerical records; the original notebook plots saved runs. Lower-rate folders \texttt{cartpole\_lr005/} and \texttt{cartpole\_lr002/} provide \texttt{run.sh} for their complete suites.

\textbf{Pong:} install its requirements and copy source/config/supervisor into a fresh directory. Set one GPU UUID in \texttt{config.json} and \texttt{supervise.sh}, and configure the supervisor's Python executable. Activate that environment:
\begin{verbatim}
export CUDA_VISIBLE_DEVICES=<single-GPU-UUID>
export CUBLAS_WORKSPACE_CONFIG=:4096:8
export OMP_NUM_THREADS=1 OPENBLAS_NUM_THREADS=1
python run_pong.py preflight
nohup setsid bash supervise.sh > nohup.log 2>&1 < /dev/null &
\end{verbatim}
Preflight runs the 15 checks; training alternates methods on one GPU and probes variance every 4,096 games. Existing checkpoints resume. Evidence: \texttt{preflight/results.json}, \texttt{runs/*/training.csv}, \texttt{runs/*/evaluation/test\_final.json} and \texttt{diagnostics/}. The README explains recovery and traces. The original suite took 2.64 hours; CartPole used CPU only.

\textbf{Report:} run \texttt{python build\_assets.py}, then compile this \LaTeX{} source. The builder audits numerical results and regenerates figures. \texttt{python combine\_submission.py} adds the shuttle PDF.

\section{Experiments with three additional months}
\textbf{CartPole:} use more training seeds and map recovery across initial positions, angles and velocities. Compare actor rate, batch size and critic-fitting frequency under matched budgets. Measure critic error and paired gradient variance throughout training to test whether critic lag explains instability.

\textbf{Pong:} keep v5, actions 2/3 and $\gamma=0.99$ while comparing frame differences, frame stacks and a convolutional policy. Measure gradient contributions by decision time. Tune rate, clipping and batch size on validation seeds; evaluate multiple training seeds on untouched 500-game tests.

\textbf{Baselines:} compare scalar, state-dependent and time-dependent predictions at identical frozen actors with independent fitting data. Use hundreds of paired games, variance-ratio confidence intervals, and separate clipped/unclipped measurements to test whether variance reduction improves control.
\end{document}
